\chapter{Fazit und Ausblick}\label{ch:conclusion}

\section{Beantwortung der Forschungsfragen}\label{sec:answers}

\paragraph{FF0 (Hauptfrage).} Wie stark aktive, messgetriebene cache-bewusste Entscheidungen ihre
passiven, statisch ausgelegten Gegenstücke übertreffen, lässt sich abschließend erst nach den
End-to-End-Läufen je Plattform beziffern (Abschnitt~\ref{sec:outlook}). Vorbereitet ist die Apparatur, die
genau diese Differenz isoliert: dieselbe achsen-konstante Struktur wird einmal aktiv (messgetriebene
Layout-, Prefetch- und Value-Ablage-Entscheidung) und einmal passiv (statisch ausgelegt) permutiert und
unter identischer Compiler- und Flag-Basis auf Hybrid-CPUs und Sapphire-Rapids gegenübergestellt.

\paragraph{FF1 (Komposition).} Die neunzehn Permutations-Achsen (Search-Algorithm,
Cache-Traversal, Mapping, Path-Compression, Node-Type, Memory-Layout, Allocator, Prefetch,
Concurrency, Serialization, Telemetry, Value-Handle, ISA, Index-Organization, I/O-Dispatch,
Migration, Filter und die beiden Queuing-Achsen) spannen einen Konfigurationsraum von weit
mehr als $10^{11}$ Permutationen auf. Stand-der-Technik-Entwürfe werden über
Algorithmus-Profil-XMLs eindeutig rekonstruiert: 30 SOTA-Profile (8 Rang-1 + 22 Rang-2/3)
zeigen, dass die Achsen-Zerlegung P01--P32 trägt. Die variadische API-Spezialisierung
(1/2/$N$ Parameter auf \texttt{std::vector}/\texttt{std::map}/Tuple) macht sowohl
algorithmus- als auch container-orientierte Muster auf einer Engine ausdrückbar. Die nicht
voll-orthogonalen Achsen-Kopplungen (Knotentyp $\times$ SIMD $\times$ Pfadkompression) werden dabei als
solche markiert und bei der Permutation explizit eingeschränkt.

\paragraph{FF2 (Mess-Methodik).} Die Messung ist in einen Produktiv-Modus
(\texttt{COMDARE\_EXPERIMENT\_MODE=OFF}, kein Overhead) und einen Experiment-Modus (der
ResultAggregator ist integraler ExecutionEngine-Member) geteilt. Profil-bewusstes
Workload-Routing wählt den Workload je Permutation aus dem Traversal-Tag und sichert so faire
Vergleiche zwischen bereichsabfrage- und punktzugriffsorientierten Entwürfen. Jede Messreihe
wird in drei Granularitäten erhoben --- Micro-, Makro- und Gesamt-Benchmarking ---, sodass
sowohl einzelne Entwurfsbestandteile als auch ganze Algorithmen vergleichbar werden. Widersprüchliche
Literaturbefunde zur optimalen Knoten-/Cache-Line-Größe (eine gegenüber bis zu sechzehn Cache-Lines) werden
so nicht aus inkompatiblen Original-Setups übernommen, sondern achsen-isoliert ceteris paribus nachgemessen.

\paragraph{FF3 (Prüfling-Vergleich).} Der konkrete Vergleich von PRT-ART gegen die 30
SOTA-Profile steht nach den ersten Mess-Läufen aus (Abschnitt~\ref{sec:outlook}). Die
Architektur ist vollständig vorbereitet: das PRT-ART-Profil, die drei
Pflicht-Messreihen-Templates, der \texttt{ExperimentDriver} mit Auto-Pickup der SOTA-Profile
und der ABI-konforme Adapter mit allen \texttt{std::map}-Operationen und Notify-Hooks. Auch die günstigste
lokale Seitendarstellung je Verzweigungspunkt (Punkt- gegenüber Kompaktseite) und ihre plattform-kalibrierten
Umschaltregeln sind Teil des PRT-ART-Profils.

\paragraph{FF4 (Compile-Time vs.\ Laufzeit).} Die Trennung ist im Build-System verankert:
plattform-spezifische Achsen (ISA, Hardware-Konstanten) werden zur Übersetzungszeit fixiert und erzeugen
provenienz-nachweisbare Binaries je Permutation (SHA256-validierte Original-Bodies,
\texttt{experiment\_compiler()}-Kennung), während Layout- und Seiten-Selektion über einen bit-codierten
Permutations-Identifier zur Laufzeit wählbar bleiben, ohne die gemeinsame Compiler- und Flag-Basis zu
verletzen.

\section{Limitierungen}\label{sec:limitations}

\begin{itemize}
  \item \textbf{Modul-Bodies sind derzeit Stubs:} die Codegenerierungs-Pipeline erzeugt
        ABI-konforme Skelett-Module je Profil; die eigentliche Such-Algorithmus-Implementierung
        je Permutation bleibt für eine Folge-Phase offen.
  \item \textbf{Hardware-Validierung ausstehend:} die Konfiguration ist lokal grün, aber der
        vollständige \texttt{ctest}-Lauf und die End-to-End-\texttt{messung\_driver}-Läufe über
        die BlockAO-Plattformen (AVX2/AVX-512, ARM NEON/SVE) stehen aus.
  \item \textbf{Ergebnis-Kapitel:} aktuell methodisch; konkrete Messungen folgen.
\end{itemize}

\section{Ausblick}\label{sec:outlook}

Kurzfristig: der End-to-End-Lauf der drei Pflicht-Messreihen auf realer Hardware, die
Generierung der Mess-Diagramme und die Vervollständigung der Modul-Bodies je Permutation.
Mittel- bis langfristig: ein GPU-Adapter (Grace-Hopper-Cluster, ZIH), die
PIM-Allokator-Integration (A16) und die formale Verifikation ausgewählter Permutations-Klassen
(StarMalloc-Stil, A13). Die Drei-Repository-Architektur erlaubt es künftigen Studierenden,
weitere Prüfling-Algorithmen analog zu PRT-ART einzubringen, ohne die
Cache-Engine-Werkzeug-Bibliothek zu modifizieren --- ein zentrales Designprinzip für die
Skalierbarkeit der Forschungsinfrastruktur. Über die hier behandelten baumartigen
Suchstrukturen hinaus ist das Anatomie-Prinzip über höherwertige Abstraktionen auf die
weiteren Strukturklassen aus Abschnitt~\ref{sec:search-classes} erweiterbar --- hash-basierte,
räumliche sowie mengen- und sequenzorientierte Strukturen ---, sodass derselbe Mess- und
Permutationsapparat ohne Neubau auf ganz neue Klassen angewendet werden kann; darin liegt ein
wesentlicher Hebel für künftigen experimentellen Fortschritt.
